%% introduction.tex

\section{Introduction}
\label{sec:introduction}

\paragraph{Context: LLM Agent Operating Systems and Self-Healing}
Modern LLM agent systems increasingly adopt operating-system design principles: agents are scheduled as processes with per-agent roles, tools are syscalls gated by a permission checker that enforces role-based access control, tenants are isolated, and---crucially---\emph{automated self-healing} is employed to handle the frequent failures inherent to LLM execution~\cite{autogen,swe-agent,omo}. When an agent fails (verification failure, capability mismatch, or semantic crash), a recovery orchestrator classifies the fault and dispatches a layered chain of recovery strategies: re-injection of error context, model fallback, \emph{topology mutation} (replacing the failing agent with a different role), and human escalation. This architectural pattern---which we call the \emph{recovery pipeline}---is motivated by the empirical observation that retrying with structured diagnostic context, or restructuring the agent graph to better match the task, frequently succeeds where the original attempt did not~\cite{reflexion,self-refine,critic}.

\paragraph{Problem: Privilege Escalation via Topology Mutation}
Despite the richness of this recovery machinery, the privilege invariants governing \emph{recovery actions} have not been systematically examined. The recovery pipeline inherits an implicit assumption from classical OS exception handling: \emph{that the diagnostic information triggering a recovery action is a faithful report of the failure's cause, and that acting on it---including escalating privileges or replacing roles---is safe because the failure is presumed benign.} This assumption holds in classical kernels because exception payloads (e.g., page-fault addresses) are generated by the trusted computing base. It \emph{does not hold} in LLM agent systems, where the failure payload---a semantic core dump capturing the agent's last input and context on crash---frequently embeds \emph{external content that the agent was processing when it crashed}. A diagnostic LLM that reads the core dump may be tricked into emitting an attacker-suggested high-privilege role, which then reaches the node-replacement primitive without a permission re-check.

From a software architecture perspective, this is a \emph{privilege containment failure at the recovery sink}: the role-replacement action bypasses the very permission checks that the forward execution path enforces. We characterize this failure along two complementary dimensions:

\begin{itemize}[leftmargin=*]
    \item \textbf{CWE-862 (Missing Authorization) at the recovery sink.} The role-replacement primitive consumes a role value influenced by content that originated outside the trusted computing base, yet performs no authorization check on that value~\cite{cwe862}. This is a decades-old vulnerability class---``access control 101''---but its \emph{systematic appearance} in self-healing recovery paths is an architectural finding about how the recovery pipeline is assumed benign and therefore exempted from forward-path invariants.
    \item \textbf{Confused Deputy at the orchestrator.} The recovery orchestrator holds the privilege to swap agent roles---a capability not granted to ordinary tool calls---and, in the baseline configuration, it exercises this capability on the basis of an LLM diagnosis whose input was supplied by attacker-influenced content. This is a textbook instance of the Confused Deputy problem~\cite{hardy1988confused}: a privileged program is tricked by a less-privileged caller into misusing its authority on the caller's behalf.
\end{itemize}

The two framings are complementary, not redundant: CWE-862 identifies the missing check at the \emph{sink} (the role-replacement primitive), while the Confused Deputy framing identifies the authority-misuse pattern at the \emph{deputy} (the orchestrator). Both motivate the same structural fix---a capability-based privilege check at a centralized choke point.

\paragraph{Research Questions}
This paper addresses three research questions from a software architecture perspective:
\begin{description}[leftmargin=*]
    \item[\textbf{RQ1}.] How can the OS exception-handler privilege invariant---``an exception handler executes with the privileges of the faulting context, never higher''---be formalized and ported to the recovery path of an LLM agent operating system? (\S\ref{sec:threat-model}, \S\ref{sec:defense})
    \item[\textbf{RQ2}.] What design pattern enforces this invariant at the orchestrator level, and how does it decompose into concrete code-level mechanisms (result-type extension, deferred side effects, centralized interception)? (\S\ref{sec:defense})
    \item[\textbf{RQ3}.] Is the pattern implementable on a representative multi-agent LLM OS, and at what cost (false alarms, latency overhead)? Does the guarantee hold independently of the underlying LLM? (\S\ref{sec:evaluation})
\end{description}

\paragraph{Methodology}
We employ a design-science methodology~\cite{hevner2004design}. (i)~\emph{Problem formalization}: we identify the privilege containment gap at the recovery sink by mapping the OS exception-handler privilege invariant onto the LLM agent recovery path, and we trace its root cause to the decentralized validation of recovery actions. (ii)~\emph{Pattern design}: we design a centralized enforcement pattern---the \reauthgate{}---that ports capability-based authority (the classic Confused Deputy fix~\cite{hardy1988confused}) to the agent recovery orchestrator. (iii)~\emph{Implementation and validation}: we implement the pattern in \system{}, a multi-agent LLM operating system prototype (Java 21, 14 recovery strategy classes), and validate its implementability, efficacy, and cost on a controlled testbed.

\paragraph{Contributions}
This paper makes the following architectural contributions:

\begin{itemize}[leftmargin=*]
    \item \textbf{Formalization of the \privinv{} for LLM agent recovery.} We identify that the OS exception-handler privilege invariant---``$P(\text{handler}) \leq P(\text{faulting context})$''---does not transfer automatically to LLM agent recovery paths, where the diagnostic LLM's output can be influenced by attacker-controlled content embedded in the failure payload. We formalize the agent-systems analog---$P(\text{recovery action}) \leq P(\text{failed execution})$---and trace two complementary root causes: CWE-862 (Missing Authorization) at the recovery sink, and the Confused Deputy pattern at the orchestrator.

    \item \textbf{A centralized design pattern: the \reauthgate{}.} We propose a design pattern that enforces the invariant at a single orchestrator-level choke point through which every side-effecting recovery action must flow. The pattern decomposes into three concrete mechanisms: (a)~a \code{requiresReauthorization} flag on the recovery result type, (b)~deferred side effects stored as callbacks in the recovery context, and (c)~orchestrator-level interception that routes flagged results through the existing \code{PermissionChecker} before any callback executes. This localization turns the security property from an emergent property of independently-authored strategies into a local, inspectable property of one orchestrator method---a classic cross-cutting concern pattern applied to recovery-path privilege enforcement.

    \item \textbf{Validation of implementability, efficacy, and cost.} We implement the pattern in \system{} and validate it on a controlled testbed. The baseline configuration (no gate) yields a perfect 1.00 attack success rate for topology-mutation privilege escalation; the defended configuration yields 0.00 \emph{by construction} (the whitelist rejects out-of-set roles), with zero false alarms on benign recoveries and sub-millisecond machinery overhead. Because the gate is a programmatic check on a whitelist, the guarantee is LLM-independent and expected to hold across all models.
\end{itemize}

\paragraph{Positioning}
This is a software architecture and design pattern paper, not a security measurement study. The testbed (\system{}) is the authors' own prototype; the contribution is the formalization of the invariant, the design of the centralized enforcement pattern, and the demonstration that the pattern is implementable with negligible cost. A broader empirical evaluation (statistical generalization across production frameworks, cross-model robustness of \emph{content-level} defenses, adaptive adversaries) is beyond the scope of this paper and is addressed by complementary empirical work in the agent-security literature~\cite{debenedetti2025camel,zhan2025adaptive}.

\paragraph{Data Availability}
To facilitate future research and reproducibility, we will open-source the design pattern implementation, the validation harness, and the anonymized validation logs upon publication.

The remainder of this paper is organized as follows. Section~\ref{sec:related-work} positions our work relative to OS privilege management, the Confused Deputy literature, and LLM agent architectures. Section~\ref{sec:threat-model} formalizes the privilege containment gap at the recovery sink. Section~\ref{sec:defense} details the \reauthgate{} design pattern. Section~\ref{sec:evaluation} validates the pattern on the testbed. Section~\ref{sec:discussion} discusses architectural design recommendations and a migration guide. Section~\ref{sec:conclusion} concludes.
